% Bagian: sets-functions-relations
% Bab: relations
% Subbagian: relations-as-sets

\documentclass[../../../include/open-logic-section]{subfiles}

\begin{document}

\olfileid[id]{sfr}{rel}{set}
\olsection{Relasi sebagai Himpunan}

\begin{explain}
Dalam \olref[sfr][set][imp]{sec}, kita menyebut beberapa himpunan penting:
$\Nat$, $\Int$, $\Rat$, $\Real$. Anda tentu mengingat beberapa relasi menarik
antara !!{element}s dari sebagian himpunan tersebut. Sebagai contoh, masing-
masing himpunan itu mempunyai \emph{relasi urutan} yang sepenuhnya baku. Ada
pula relasi \emph{identik dengan}, yang berlaku antara setiap objek dan
dirinya sendiri, tetapi tidak dengan objek lain. Kita akan menjumpai jauh
lebih banyak relasi menarik, dan bahkan lebih banyak lagi relasi yang mungkin.
Namun, sebelum meninjaunya, kita akan mulai dengan menunjukkan bahwa relasi
dapat dipandang sebagai suatu jenis himpunan yang khusus.

Untuk itu, ingatlah dua hal dari \olref[sfr][set][pai]{sec}. Pertama, ingat
gagasan \emph{pasangan terurut}: jika diberikan $a$ dan $b$, kita dapat
membentuk~$\tuple{a, b}$. Yang penting, urutan elemen memang \emph{berpengaruh}
di sini. Jadi, jika $a \neq b$, maka $\tuple{a, b} \neq \tuple{b, a}$.
(Bandingkan dengan pasangan tak terurut, yakni himpunan beranggota $2$, yang
memenuhi $\{a, b\}=\{b, a\}$.) Kedua, ingat gagasan \emph{hasil kali
Kartesius}: jika $A$ dan $B$ adalah himpunan, kita dapat membentuk~$A \times
B$, yaitu himpunan semua pasangan $\tuple{x, y}$ dengan $x \in A$ dan $y \in
B$. Secara khusus, $A^{2}= A \times A$ adalah himpunan semua pasangan terurut
dari~$A$.

Sekarang kita akan meninjau suatu relasi tertentu pada sebuah himpunan, yaitu
relasi~$<$ pada himpunan bilangan asli~$\Nat$. Perhatikan himpunan semua
pasangan bilangan $\tuple{n, m}$ dengan $n<m$, yakni
\[
  R=\Setabs{\tuple{n, m}}{n, m \in \Nat \text{ dan } n<m}.
\]
Terdapat hubungan erat antara $n$ yang lebih kecil daripada $m$ dan pasangan
$\tuple{n, m}$ yang menjadi anggota $R$, yaitu
\[
      n<m\text{ jika dan hanya jika }\tuple{n, m} \in R.
\]
Sesungguhnya, tanpa kehilangan informasi apa pun, kita dapat memandang
himpunan $R$ \emph{sebagai} relasi~$<$ pada $\Nat$.

Dengan cara yang sama, kita dapat membentuk himpunan bagian dari $\Nat^{2}$
untuk setiap relasi antara bilangan. Sebaliknya, jika diberikan suatu himpunan
pasangan bilangan $S \subseteq \Nat^{2}$, terdapat relasi yang bersesuaian
antara bilangan, yaitu relasi yang berlaku antara $n$ dan $m$ jika dan hanya jika
$\tuple{n, m} \in S$. Hal ini membenarkan definisi berikut.
\end{explain}

\begin{defn}[Relasi biner]
Suatu \emph{relasi biner} pada himpunan $A$ adalah himpunan bagian
dari~$A^{2}$. Jika $R \subseteq A^{2}$ merupakan relasi biner pada~$A$ dan
$x, y \in A$, kita kadang-kadang menulis $Rxy$ (atau $xRy$) sebagai ganti
$\tuple{x, y} \in R$.
\end{defn}

\begin{ex}
  \ollabel{relations}
Himpunan $\Nat^{2}$ yang terdiri atas pasangan bilangan asli dapat dicantumkan
dalam matriks dua dimensi seperti berikut.
\[
  \begin{array}{ccccc}
  \mathbf{\tuple{ 0,0 }} & \tuple{ 0,1 } &
    \tuple{ 0,2 } & \tuple{ 0,3 } & \ldots\\
  \tuple{ 1,0 } & \mathbf{\tuple{ 1,1 }} &
    \tuple{ 1,2 } & \tuple{ 1,3 } & \ldots\\
  \tuple{ 2,0 } & \tuple{ 2,1 } &
    \mathbf{\tuple{ 2,2 }} & \tuple{ 2,3 } & \ldots\\
  \tuple{ 3,0 } & \tuple{ 3,1 } & \tuple{ 3,2 } &
    \mathbf{\tuple{ 3,3 }} & \ldots\\
  \vdots & \vdots & \vdots & \vdots & \mathbf{\ddots}
  \end{array}
\]
Di sini kita mencetak diagonal dengan huruf tebal karena himpunan bagian dari
$\Nat^2$ yang terdiri atas pasangan pada diagonal, yakni
\[
  \{\tuple{0,0 }, \tuple{ 1,1 }, \tuple{ 2,2 }, \dots\},
  \]
merupakan \emph{relasi identitas pada}~$\Nat$. (Karena relasi identitas sering
digunakan, mari kita definisikan $\Id{A}=\Setabs{\tuple{ x,x }}{x \in A}$
untuk setiap himpunan $A$.) Himpunan bagian yang terdiri atas
semua pasangan di atas diagonal, yakni
\[
  L = \{\tuple{ 0,1 },\tuple{ 0,2 },\ldots,\tuple{ 1,2 },
  \tuple{ 1,3 }, \dots, \tuple{ 2,3 }, \tuple{ 2,4 },\ldots\},
\]
merupakan relasi \emph{lebih kecil daripada}, yakni $Lnm$ jika dan hanya jika
$n<m$. Himpunan bagian yang terdiri atas pasangan di bawah diagonal, yakni
\[
  G=\{ \tuple{ 1,0 },\tuple{ 2,0 },\tuple{
    2,1 }, \tuple{ 3,0 },\tuple{ 3,1 },\tuple{ 3,2 }, \dots\},
\]
merupakan relasi \emph{lebih besar daripada}, yakni $Gnm$ jika dan hanya jika
$n>m$. Gabungan $L$ dengan $\Id{\Nat}$, yang dapat kita sebut
$K=L\cup \Id{\Nat}$, merupakan
relasi \emph{lebih kecil daripada atau sama dengan}: $Knm$ jika dan hanya jika
$n \le m$. Demikian pula, $H=G \cup \Id{\Nat}$ merupakan relasi \emph{lebih besar
daripada atau sama dengan}. Relasi $L$, $G$, $K$, dan $H$ merupakan jenis
relasi khusus yang disebut \emph{urutan}. Relasi $L$ dan $G$ memiliki sifat
bahwa tidak ada bilangan yang berelasi $L$ atau $G$ dengan dirinya sendiri
(yakni, untuk setiap $n$, $Lnn$ maupun $Gnn$ tidak berlaku). Relasi dengan
sifat ini disebut \emph{irefleksif}; jika relasi tersebut juga merupakan
urutan, relasi itu disebut \emph{urutan ketat}.
\end{ex}

\begin{explain}
Meskipun urutan dan identitas merupakan relasi yang penting dan alami, perlu
ditekankan bahwa menurut definisi kita, \emph{setiap} himpunan bagian dari
$A^{2}$ merupakan relasi pada~$A$, betapapun tidak alami atau dibuat-buatnya
relasi tersebut. Secara khusus, $\emptyset$ merupakan relasi pada setiap
himpunan (\emph{relasi kosong}, yang tidak berlaku pada pasangan elemen mana
pun), sedangkan $A^{2}$ sendiri juga merupakan relasi pada~$A$ (relasi yang
berlaku pada setiap pasangan), yang disebut \emph{relasi universal}. Bahkan
sesuatu seperti $E=\Setabs{\tuple{n, m}}{n>5 \text{ atau } m \times n \ge 34}$
pun tergolong relasi.
\end{explain}

\begin{prob}
  Tuliskan semua !!{element}s relasi $\subseteq$ pada himpunan
  $\Pow{\{a, b, c\}}$.
\end{prob}

\end{document}
